\section{Envelope and solver details}
\label{supp:details}

Two implementation notes referenced from the main text: how each family spends
the envelope's advancing index, and how to read the two below-unity rows of the
ablation.

\subsection{Every family has a currency}

The envelope view advances every layer's index together by $u$ and averages over
$u\in[0,1)$, and each family spends that advance in its own coin. The level-set
families shift a phase. The radial pencil turns: advancing its index by $u$
rotates the fan by $u\pi/N$, which is how a family whose index is an angle joins
an average built for scalars. A lattice indexes its members by a \emph{pair} of
integers, so it translates along its two generators --- and here the diagonal is
exactly wrong. Stepping both generators in lockstep keeps their two indices
correlated, and a lattice then beats against itself: a fine diagonal hatch at the
carrier's own scale that no amount of contrast can be read through. So we take
the $i$th tap along the first generator to $u_i$ and along the second to
$\{i\varphi\}$, $\varphi$ the golden ratio --- the Fibonacci lattice, a rank-1
rule on the two-torus~\cite{Sloan1994Lattice, Keller2013Quasi}. It decorrelates a
lattice from itself while leaving it in step with every other layer's matching
generator, which is where the fringes an author wants actually are.

\subsection{A pivot, not a gain}

The view exposes a contrast, because a fringe field is a small excursion about
the stack's mean coverage and the interesting part is the excursion. Expanding
about $0$ would ride the mean up and down as the fringes move, so the expansion
is about the coverage the stack would average over its own periods: $2h/s$ per
family, composited. Counting a lattice's families is not optional. A hex lattice
draws three, and calling its coverage that of one puts the pivot far brighter
than the picture it is expanding about, which at any real contrast drives the
whole frame black --- which is how we found the bug.

\subsection{Reading the ablation's negative rows}

Removing the carried rotation or the accept-early-exit makes the dense ablation
scenes $8$--$19\%$ \emph{faster}. Removing the accept exit is faster because in a
heavily inked region almost every lane would exit anyway, and the branch only
adds divergence: with $32$ lanes in lockstep, one slow lane sets the pace, so an
early exit that helps $31$ of them helps nobody. Both mechanisms stay because
they are what make sparse and mid-density scenes cheap, and because a solver
tuned only for saturated fields is tuned for the case where the image carries
least information. The same wavefront-divergence argument says why narrowing the
interval, rather than exiting the loop sooner, is the productive lever in this
problem: the interval is computed per pixel before the loop and shrinks the work
of \emph{every} lane, while an exit condition only helps whichever lane is
already done.
